% GENERATED FILE. Edit canonical lessons, then run npm run generate:books.
% canonical-source-hash: fnv1a64:f92c4916ec153153
% canonical-lessons: GU-C248-bethak, GU-C248-arji, GU-C248-sahi-karvi, GU-C248-sopvu, GU-C248-nodhvu, GU-R248-review-243-248

\chapter{Meetings and applications}
\label{ch:gu-meetings-and-applications}
\hlchaptermodality{248}

\begin{chapteropening}
\textbf{By the end of this chapter:} \emph{I can talk about meetings, applications and signing papers.}

Complete the last lesson of chapter 248: I can talk about meetings, applications and signing papers.
\end{chapteropening}

\section[\texorpdfstring{\gu{બેઠક} (beṭhak)}{બેઠક (beṭhak)}]{\gu{બેઠક} (beṭhak) — meeting; sitting room}

\label{lesson:GU-C248-bethak}



\begin{quote}
Before the new one: say the Gujarati for to earn, then the Gujarati for hardworking.
\end{quote}

\subsection*{You'll want to know: \gu{બેઠક}}
\textbf{\gu{બેઠક}} — \emph{beṭhak} — \textquotedblleft{}meeting; sitting room\textquotedblright{}.

\begin{quote}
\textbf{\gu{આજે} \gu{બપોરે} \gu{બેઠક} \gu{છે}.} — \emph{āje bapore beṭhak che} — \textquotedblleft{}There is a meeting this afternoon.\textquotedblright{}
\end{quote}

A feminine noun.

\begin{cousinweb}
\textbf{\gu{બેઠક}} (\emph{beṭhak}), a sitting together, is built on \emph{beṭh-}, a stem of the taught \emph{besvũ}, \textquotedblleft{}to sit\textquotedblright{}, with \emph{-ak}. CDIAL derives \emph{besvũ} and \emph{beṭho} from Sanskrit \emph{úpaviśati}, \textquotedblleft{}sits down\textquotedblright{} (CDIAL 2245).
\end{cousinweb}

\begin{grammarlens}[title={What you've built}]
One more word for everyday life at A2.
\end{grammarlens}

\subsection*{Your turn}
\begin{itemize}
\raggedright
  \item \emph{Say it:} \emph{beṭhak}
  \item \emph{Say it:} \emph{beṭhak}, once more
  \item \emph{Say it:} \emph{āje bapore beṭhak che} — There is a meeting this afternoon.
  \item \emph{Recall:} say the Gujarati for to earn, then the Gujarati for hardworking, then say \emph{beṭhak} again
\end{itemize}

\begin{tcolorbox}[breakable,colback=teal!4,colframe=teal!35!black,title={Before you move on}]
What does \gu{બેઠક} mean? (\textquotedblleft{}Meeting; sitting room\textquotedblright{}.) Say it once more.
\end{tcolorbox}
\section[\texorpdfstring{\gu{અરજી} (arjī)}{અરજી (arjī)}]{\gu{અરજી} (arjī) — application, petition}

\label{lesson:GU-C248-arji}



\begin{quote}
Before the new one: say the Gujarati for hardworking, then the Gujarati for meeting.
\end{quote}

\subsection*{You'll want to know: \gu{અરજી}}
\textbf{\gu{અરજી}} — \emph{arjī} — \textquotedblleft{}application, petition\textquotedblright{}.

\begin{quote}
\textbf{\gu{તે} \gu{નોકરી} \gu{માટે} \gu{અરજી} \gu{લખે} \gu{છે}.} — \emph{te nokrī māṭe arjī lakhe che} — \textquotedblleft{}She writes an application for a job.\textquotedblright{}
\end{quote}

A feminine noun.

\begin{cousinweb}
\textbf{\gu{અરજી}} (\emph{arjī}) is a loan usually traced to a Persian \emph{ʿarzī}, \textquotedblleft{}a petition\textquotedblright{}, built on Arabic \emph{ʿarḍ}, \textquotedblleft{}a presenting, a showing\textquotedblright{}. Its Arabic root is \emph{ʿ-r-ḍ} (Lane).
\end{cousinweb}

\begin{grammarlens}[title={What you've built}]
One more word for everyday life at A2.
\end{grammarlens}

\subsection*{Your turn}
\begin{itemize}
\raggedright
  \item \emph{Say it:} \emph{arjī}
  \item \emph{Say it:} \emph{arjī}, once more
  \item \emph{Say it:} \emph{te nokrī māṭe arjī lakhe che} — She writes an application for a job.
  \item \emph{Recall:} say the Gujarati for hardworking, then the Gujarati for meeting, then say \emph{arjī} again
\end{itemize}

\begin{tcolorbox}[breakable,colback=teal!4,colframe=teal!35!black,title={Before you move on}]
What does \gu{અરજી} mean? (\textquotedblleft{}Application, petition\textquotedblright{}.) Say it once more.
\end{tcolorbox}
\section[\texorpdfstring{\gu{સહી} \gu{કરવી} (sahī karvī)}{સહી કરવી (sahī karvī)}]{\gu{સહી} \gu{કરવી} (sahī karvī) — to sign}

\label{lesson:GU-C248-sahi-karvi}



\begin{quote}
Before the new one: say the Gujarati for meeting, then the Gujarati for application.
\end{quote}

\subsection*{You'll want to know: \gu{સહી} \gu{કરવી}}
\textbf{\gu{સહી} \gu{કરવી}} — \emph{sahī karvī} — \textquotedblleft{}to sign\textquotedblright{}.

\begin{quote}
\textbf{\gu{સાહેબ} \gu{સાંજે} \gu{બધા} \gu{કાગળમાં} \gu{સહી} \gu{કરે} \gu{છે}.} — \emph{sāheb sā̃je badhā kāgaḷmā̃ sahī kare che} — \textquotedblleft{}The boss signs all the papers in the evening.\textquotedblright{}
\end{quote}

\textbf{\gu{સહી}}, \emph{sahī}, alone means a signature.

\begin{cousinweb}
\textbf{\gu{સહી} \gu{કરવી}} (\emph{sahī karvī}) is built from \emph{sahī}, \textquotedblleft{}signature\textquotedblright{}, and \emph{karvī}, \textquotedblleft{}to do\textquotedblright{}. \emph{Sahī} is usually traced to Arabic \emph{ṣaḥīḥ}, \textquotedblleft{}sound, correct\textquotedblright{}, through Persian; its Arabic root is \emph{ṣ-ḥ-ḥ} (Lane).
\end{cousinweb}

\begin{grammarlens}[title={What you've built}]
One more word for everyday life at A2.
\end{grammarlens}

\subsection*{Your turn}
\begin{itemize}
\raggedright
  \item \emph{Say it:} \emph{sahī karvī}
  \item \emph{Say it:} \emph{sahī karvī}, once more
  \item \emph{Say it:} \emph{sāheb sā̃je badhā kāgaḷmā̃ sahī kare che} — The boss signs all the papers in the evening.
  \item \emph{Recall:} say the Gujarati for meeting, then the Gujarati for application, then say \emph{sahī karvī} again
\end{itemize}

\begin{tcolorbox}[breakable,colback=teal!4,colframe=teal!35!black,title={Before you move on}]
What does \gu{સહી} \gu{કરવી} mean? (\textquotedblleft{}To sign\textquotedblright{}.) Say it once more.
\end{tcolorbox}
\section[\texorpdfstring{\gu{સોંપવું} (sõpvũ)}{સોંપવું (sõpvũ)}]{\gu{સોંપવું} (sõpvũ) — to hand over, to entrust}

\label{lesson:GU-C248-sopvu}



\begin{quote}
Before the new one: say the Gujarati for application, then the Gujarati for to sign.
\end{quote}

\subsection*{You'll want to know: \gu{સોંપવું}}
\textbf{\gu{સોંપવું}} — \emph{sõpvũ} — \textquotedblleft{}to hand over, to entrust\textquotedblright{}.

\begin{quote}
\textbf{\gu{સાહેબ} \gu{મને} \gu{નવું} \gu{કામ} \gu{સોંપે} \gu{છે}.} — \emph{sāheb mane navũ kām sõpe che} — \textquotedblleft{}The boss hands me a new task.\textquotedblright{}
\end{quote}

\begin{cousinweb}
\textbf{\gu{સોંપવું}} (\emph{sõpvũ}) is a tadbhava, a word inherited from Sanskrit and reshaped by sound change. It comes from \emph{sámarpayati}, \textquotedblleft{}delivers over\textquotedblright{}; CDIAL glosses the Gujarati verb \textquotedblleft{}hands over\textquotedblright{} (CDIAL 13192).
\end{cousinweb}

\begin{grammarlens}[title={What you've built}]
One more word for everyday life at A2.
\end{grammarlens}

\subsection*{Your turn}
\begin{itemize}
\raggedright
  \item \emph{Say it:} \emph{sõpvũ}
  \item \emph{Say it:} \emph{sõpvũ}, once more
  \item \emph{Say it:} \emph{sāheb mane navũ kām sõpe che} — The boss hands me a new task.
  \item \emph{Recall:} say the Gujarati for application, then the Gujarati for to sign, then say \emph{sõpvũ} again
\end{itemize}

\begin{tcolorbox}[breakable,colback=teal!4,colframe=teal!35!black,title={Before you move on}]
What does \gu{સોંપવું} mean? (\textquotedblleft{}To hand over, to entrust\textquotedblright{}.) Say it once more.
\end{tcolorbox}
\section[\texorpdfstring{\gu{નોંધવું} (nõdhvũ)}{નોંધવું (nõdhvũ)}]{\gu{નોંધવું} (nõdhvũ) — to note down}

\label{lesson:GU-C248-nodhvu}



\begin{quote}
Before the new one: say the Gujarati for to sign, then the Gujarati for to hand over.
\end{quote}

\subsection*{You'll want to know: \gu{નોંધવું}}
\textbf{\gu{નોંધવું}} — \emph{nõdhvũ} — \textquotedblleft{}to note down\textquotedblright{}.

\begin{quote}
\textbf{\gu{હું} \gu{બેઠકમાં} \gu{બધું} \gu{નોંધું} \gu{છું}.} — \emph{hũ beṭhakmā̃ badhũ nõdhũ chũ} — \textquotedblleft{}I note down everything in the meeting.\textquotedblright{}
\end{quote}

Its causative \textbf{\gu{નોંધાવવું}} means to register.

\begin{cousinweb}
\textbf{\gu{નોંધવું}} (\emph{nõdhvũ}) is a tadbhava. It comes from \emph{*nibandhati}, beside Sanskrit \emph{níbadhnāti}, \textquotedblleft{}binds on; writes down\textquotedblright{}; CDIAL glosses the Gujarati verb \textquotedblleft{}to note down, enter in a book\textquotedblright{} (CDIAL 7223).
\end{cousinweb}

\begin{grammarlens}[title={What you've built}]
One more word for everyday life at A2.
\end{grammarlens}

\subsection*{Your turn}
\begin{itemize}
\raggedright
  \item \emph{Say it:} \emph{nõdhvũ}
  \item \emph{Say it:} \emph{nõdhvũ}, once more
  \item \emph{Say it:} \emph{hũ beṭhakmā̃ badhũ nõdhũ chũ} — I note down everything in the meeting.
  \item \emph{Recall:} say the Gujarati for to sign, then the Gujarati for to hand over, then say \emph{nõdhvũ} again
\end{itemize}

\begin{tcolorbox}[breakable,colback=teal!4,colframe=teal!35!black,title={Before you move on}]
What does \gu{નોંધવું} mean? (\textquotedblleft{}To note down\textquotedblright{}.) Say it once more.
\end{tcolorbox}
\section[(dialogue)]{Work, school and study}

\label{lesson:GU-R248-review-243-248}



\begin{quote}
Say the Gujarati for lesson and to learn by heart.
\end{quote}

\subsection*{Your turn}
Say these aloud:

\begin{itemize}
\raggedright
  \item lesson, to learn by heart, schoolbook, school bag, present
  \item mathematics, geography, essay, to explain, to solve
  \item homework, spelling, to rub out, eraser, to pronounce
  \item result, marks, to fail an exam, studies, to work hard
  \item boss, employee, business, to earn, hardworking
  \item meeting, application, to sign, to hand over, to note down
\end{itemize}

\begin{tcolorbox}[breakable,colback=teal!4,colframe=teal!35!black,title={Before you move on}]
Lesson? (\textbf{\gu{પાઠ}}.) To learn by heart? (\textbf{\gu{ગોખવું}}.) Schoolbook? (\textbf{\gu{ચોપડી}}.) School bag? (\textbf{\gu{દફતર}}.) Present? (\textbf{\gu{હાજર}}.) Mathematics? (\textbf{\gu{ગણિત}}.) Geography? (\textbf{\gu{ભૂગોળ}}.) Essay? (\textbf{\gu{નિબંધ}}.) To explain? (\textbf{\gu{સમજાવવું}}.) To solve? (\textbf{\gu{ઉકેલવું}}.) Homework? (\textbf{\gu{ગૃહકાર્ય}}.) Spelling? (\textbf{\gu{જોડણી}}.) To rub out? (\textbf{\gu{ભૂંસવું}}.) Eraser? (\textbf{\gu{રબર}}.) To pronounce? (\textbf{\gu{ઉચ્ચારવું}}.) Result? (\textbf{\gu{પરિણામ}}.) Marks? (\textbf{\gu{ગુણ}}.) To fail an exam? (\textbf{\gu{નાપાસ} \gu{થવું}}.) Studies? (\textbf{\gu{અભ્યાસ}}.) To work hard? (\textbf{\gu{મહેનત} \gu{કરવી}}.) Boss? (\textbf{\gu{સાહેબ}}.) Employee? (\textbf{\gu{કર્મચારી}}.) Business? (\textbf{\gu{ધંધો}}.) To earn? (\textbf{\gu{કમાવું}}.) Hardworking? (\textbf{\gu{મહેનતુ}}.) Meeting? (\textbf{\gu{બેઠક}}.) Application? (\textbf{\gu{અરજી}}.) To sign? (\textbf{\gu{સહી} \gu{કરવી}}.) To hand over? (\textbf{\gu{સોંપવું}}.) To note down? (\textbf{\gu{નોંધવું}}.)
\end{tcolorbox}
